\newpage
\section{Charting Table Template with Example}
\label{sec:appendix-scr-template}

% \label{sec:appendix-filled-table}
\begin{landscape}
\begin{table}[p]
  \centering
  \caption{In addition to summary information, the charting table template shows the information extracted from the papers included in the analysis. The example entry shows extracted and labelled information for a paper by \citet{vallgaarda2014dress}.}
  \label{tab:chartingtable}

  \appendixtablefont
  \verytighttable

  \begin{tabularx}{\linewidth}{@{}
    p{.16\linewidth}
    p{.12\linewidth}
    p{.10\linewidth}
    p{.13\linewidth}
    p{.20\linewidth}
    p{.12\linewidth}
    X
  @{}}
    \toprule
    Research motivations &
    Methodologies &
    Type of smart building &
    Interactive technologies &
    Human experience &
    Measures and instruments &
    Visions for future research \\
    \midrule
        \emph{``What does it entail to be embraced by a space that responds to your actions?''} & Conceptual design and prototyping of a responsive space, using modern dance as a method to explore interaction qualities. &  Interactive installation & The Dress Room - white cube that responds to the body movements over the floor & \emph{``What does it entail to be embraced by a space that responds to your actions?''} & Sensors in the floor detecting body movements, Pneumatic pistons to move the walls, Arduino Mega boards for control" "White textile tent suspended within a steel frame & Future architectural spaces will incorporate responsive elements to enhance interaction and experience \\
        \emph{``What kind of relations can we create between the active body and the active space?''} & & & &\emph{``What kind of relations can we create between the active body and the active space?''} & \emph{DANCE AS RESEARCH METHOD - Observations and interviews} & \emph{``The Dress Room is a responsive space. It moves. It adapts. The floor senses where you are. The room responds. It follows you. It stops. It collapses then expands. You are enclosed. It is a dress you wear. Sometimes it fits, sometimes it misbehaves. It invites you to move. To explore. To dance. The Dress Room blurs the boundaries between architecture and clothes.''} \\
        \emph{``What qualities does the responsivity have for creating certain experiences of a space?''} & & & & \emph{``What qualities does the responsivity have for creating certain experiences of a space?''} & & \emph{``Thus, with the Dress Room as a simple responsive space explored through modern dance I have begun to outline some of the experiential qualities of being embraced by a space that responds to our bodily motions. My explorations have indicated that responsive spaces hold the potential to make us form a kind of symbiosis with the space they are in – that the temporal form of the space’s responsivity enables us to design for a sense of intimacy. They further indicated that our actions and the space’s responses become interchanged and while the body and the room moves together they are constantly motivated/actuated by one another. The mutual responses of the body and the space are integrated but not the same. They are in a cause-and-effect relationship that is constantly evolving.''} \\
        \emph{``Through the Dress Room, I begin to explore the qualities of responsive spaces and embodied interaction.''} & & & & \emph{``. This seems to indicate that exposing our body to an unfamiliar responsive environment may be a bit daunting at first it may be something we need to overcome.''} & & \\
        \emph{``Indeed, what I am interested in with this study is to explore the kind of relations we can create between the active body and the active space. I see explorations like this as a step towards understanding what experiences and new types of interaction these responsive spaces can foster''} & & & & \emph{``The Dress Room’s ability to motivate Nana to move towards certain areas of the space as well as to move her body in specific ways indicate that responsive spaces can have a significant qualities when it comes to motion motivation.''} & &  \\
        \emph{``What I want to explore here is the experience of the responsivity from a fully embodied interaction perspective. What experiential qualities does it have? What possibly new ways of inhabiting space it can encourage? What potential new types of interaction can we develop from this?''} & & & & \emph{``Thus, with the Dress Room as a simple responsive space explored through modern dance I have begun to outline some of the experiential qualities of being embraced by a space that responds to our bodily motions. My explorations have indicated that responsive spaces hold the potential to make us form a kind of symbiosis with the space they are in – that the temporal form of the space’s responsivity enables us to design for a sense of intimacy. They further indicated that our actions and the space’s responses become interchanged and while the body and the room moves together they are constantly motivated/actuated by one another. The mutual responses of the body and the space are integrated but not the same. They are in a cause-and-effect relationship that is constantly evolving.''} & & \\
        & & & & \emph{``The Dress Room is a white cube that responds to the body’s movements over the floor. The walls move, the room collapses or expands.''} & &\\
        & & & & \emph{``The outcome suggests that interacting with responsive environments can help create a sense of intimacy as well as motivate our motions within the space.''} & & \\
        & & & & \emph{``The Dress Room is a white cube that responds to the body movements over the floor. It is made from a white square textile tent suspended within a steel frame measuring 5x5x5 meter. The suspension allows for the entire room to move more than half a meter in each direction. The sensors in the floor detect where a body is and the room responds either by moving with or away from that body depending on the setup.''} & & \\
        & & & & \emph{``Through a dancer’s experiences we begin to understand what experiential qualities, and thus what possibilities for interaction design, this kind of responsiveness affords.''} & &\\
        & & & & \emph{``The Dress Room was developed as an aesthetic exploration and is deliberately abstract in its expression. It has no function other than as a responsive enclosure.''} & &\\
        & & & & \emph{``The Dress Room is an abstract installation made in an attempt to focus the experience around the responsivity. It is conceived as a space you take on like a dress.''} & & \\
        & & & & \emph{``the Dress Room is about letting the white garment create a space in co-production with the movements of the body.''} & & \\
        & & & & \emph{``First Type Responsive Form: The Room as a Dress We started out letting the room respond as a dress in the sense that the walls would follow the body within. The room follows you around. It makes you feel that you wear a room. You feel in Control. You belong together. You cannot escape – the door keeps evading you.''} & & \\
        & & & & \emph{``Second Type of Responsive Form: The Vertical Response The second type of responsive form was designed to explore the vertical dimensions of the room. We had deliberately created the room with a high ceiling allowing us to collapse it without coming in contact with the person within''} & & \\
        & & & & \emph{``Still, what does it mean to inhabit an abstract space like the Dress Room? It is abstract like art yet it demands more than an onlooker to be understood. It demands active participation. Thus, I have chosen a method that meets this space at its level of abstraction while still enabling active bodily engagement. I have chosen to use modern dance performed by a professional dancer whose experiences I have sought to understand through observations and interviews.''} & & \\
        & & & & \emph{``Overall Nana described the Dress Room as a constant source of inspiration.  She explained that even when she repeated a series of movements within the same responsive form and the response from the pistons was the same the fabric would ripple differently each time creating an overall unique expression.''} & & \\
        & & & & \emph{``She explained that after a while she ceased to focus on how her actions influenced specific responses, and instead she eased into a dance where her motions became responses as well as actuators.''} & & \\
        & & & & \emph{``Lastly, Nana explained while demonstrating through body language how certain responses in the Dress Room would inspire her to do certain movements.''} & & \\
        & & & & \emph{``The two experiential qualities that Nana experienced in the responsive room can be articulated as intimacy and motivated motion.''} & & \\
        & & & & \emph{``With the Dress Room, I sat out to explore embodied interaction with responsive spaces.''} & & \\
        & & & & \emph{``Thus, with the Dress Room as a simple responsive space explored through modern dance I have begun to outline some of the experiential qualities of being embraced by a space that responds to our bodily motions. My explorations have indicated that responsive spaces hold the potential to make us form a kind of symbiosis with the space they are in – that the temporal form of the space’s responsivity enables us to design for a sense of intimacy. They further indicated that our actions and the space’s responses become interchanged and while the body and the room moves together they are constantly motivated/actuated by one another. The mutual responses of the body and the space are integrated but not the same. They are in a cause-and-effect relationship that is constantly evolving.''} & & \\
         \bottomrule
  \end{tabularx}
\end{table}
\end{landscape}

\newpage 

\section{General Characteristics of the Corpus Papers}

\begin{widefigure}[tbp]
  \centering

  \begin{subfigure}{.48\linewidth}
    \centering
    \includesvg[width=\linewidth]{images/findings/descriptive/years.svg}
    \caption{Papers by publication year.}
    \label{fig:overview-years}
  \end{subfigure}
  \hfill
  \begin{subfigure}{.48\linewidth}
    \centering
    \includesvg[width=\linewidth]{images/findings/descriptive/type-of-smart-building.svg}
    \caption{Papers by type of smart building.}
    \label{fig:overview-building}
  \end{subfigure}

  \vspace{.7\baselineskip}

  \begin{subfigure}{.48\linewidth}
    \centering
    \includesvg[width=\linewidth]{images/findings/descriptive/type-of-publication.svg}
    \caption{Papers by type of publication.}
    \label{fig:overview-publication}
  \end{subfigure}
  \hfill
  \begin{subfigure}{.48\linewidth}
    \centering
    \includesvg[width=\linewidth]{images/findings/descriptive/conferences.svg}
    \caption{Papers by publication venue.}
    \label{fig:overview-venue}
  \end{subfigure}

  \caption{An overview of the characteristics of the papers included in the scoping review: publication year, type of smart building, publication type, and publication venue.}
  \label{fig:findings-overview}
\end{widefigure}


\newpage 

\section{A Typology of Technological Interventions}
\begin{widetable}[tbp]
  \caption{A typology of technological interventions in smart buildings, including descriptions, number of papers, prototypical examples, and the three most frequently associated targeted experiences.}
  \label{tab:technology-types}

  \appendixwidefont
  \tighttable

  \begin{tabularx}{\linewidth}{@{}
    p{.16\linewidth}
    Y
    c
    Y
    p{.18\linewidth}
  @{}}
    \toprule
    Type of intervention &
    Description &
    N &
    Prototypical examples &
    Associated targeted experiences \\
    \midrule

    Audio and Lighting &
    Sound and light as key design materials, often synchronised to evoke emotional responses, guide behaviour, or modify spatial perception, creating multisensory and immersive environments. &
    27 &
    Mediated Atmospheres is a modular workspace using lighting, video projection, and sound to transform the appearance of the workspace \cite{richer2020exploring}. Trivet Fields uses mini-speakers to create soundscapes to aesthetically delineate space \cite{jakovich2007sonictecture}. &
    Aesthetic Experiences; Emotional Experiences; Comfort \\

    Augmented and Virtual Reality &
    Immersive experiences blending physical and digital spaces, enhancing spatial awareness and interaction in novel ways. &
    17 &
    Weightless Walls uses user positions and noise-cancelling headphones to display virtual walls for auditory privacy \cite{takeuchi2010weightless}. VR Living Rooms use virtual simulations of indoor environments to measure psychological responses and affective states \cite{tawil2021living}. &
    Explorative and/or Future and Speculative; Sense of Place-making; Emotional Experiences \\

    Environmental Sensing &
    IoT-connected sensors monitor environmental conditions such as temperature, light, air quality, and noise, either adjusting parameters automatically or providing real-time information. &
    23 &
    Environmental sensors were deployed to measure temperature, humidity, and light, with data visualised for occupants and building managers \cite{clear2017d}. Hilo-box uses CO2 sensors and machine learning to forecast indoor air quality during office meetings \cite{zhong2021complexity}. &
    Comfort; Democratic Experiences; Explorative and/or Future and Speculative \\

    Feedback Systems &
    Participatory sensing systems, mobile apps, or ambient cues that allow users to collect data, adjust environments, or receive guidance. &
    11 &
    Follow-the-Lights uses ambient lighting, pressure mats, and abstract visualisations to influence stair and elevator usage \cite{rogers2010ambient}. A building Twitter bot engages occupants around preferences and experiences \cite{mitchell2020no}. &
    Emotional Experiences; Democratic Experiences; Inclusion Experiences \\

    Interactive Displays &
    Displays integrated with cameras, motion tracking, or simulations that adapt to user interactions and create dynamic environments. &
    45 &
    An interactive window display uses large screens, RGB/depth cameras, and temperature sensors to display context-related scenery \cite{wang2015intelligent}. WaveWindow uses projectors, microphones, and cameras for interaction on a retail window display \cite{perry2010wavewindow}. &
    Explorative and/or Future and Speculative; Experience of Relatedness; Emotional Experiences \\

    Physiological Sensing &
    Biofeedback devices and wearables monitor bodily responses, with some interventions adapting environmental conditions in real time. &
    10 &
    EEG, heart rate, skin conductance, and pulse volume measurements were used to understand neurophysiological responses to smart home systems \cite{angioletti2020neurophysiological}. &
    Comfort; Emotional Experiences; Inclusion Experiences \\

    Shape-Changing and Kinetic Interfaces &
    Dynamic physical changes in environments through adaptive furniture, walls, or materials that respond to user presence and activity. &
    22 &
    TidalSpace creates an interactive spatial boundary for work-from-home parents \cite{huynh2022tidal}. The shape-changing work desk uses digital tabletops and motion tracking for proxemic transitions in informal meetings \cite{gronbaek2017proxemic}. &
    Bodily Experiences; Explorative and/or Future and Speculative; Sense of Place-making \\

    \bottomrule
  \end{tabularx}
\end{widetable}


\newpage

\section{Detailed Breakdown of Research Trends}
\label{sec:appendix-breakdown}

\begin{figure}[!htbp]
    \centering
    \includesvg[width=\linewidth]{images/findings/appendix/trends-publications-appendix.svg}
    \caption{Research trends shows all publications targeting human experiences in smart buildings over time.}
    \label{fig:appendix-rq-trends}
\end{figure}

\newpage

\section{An Overview of Scoping Review Findings}
\label{sec:appendix-unfolded-sankey}
\begin{widefigure}[tbp]
  \centering

  \begin{subfigure}{.5\linewidth}
    \centering
    \includesvg[width=\linewidth]{images/findings/appendix/unfolded-themes-mechanisms.svg}
    \caption{Human experiences and associated design mechanisms.}
    \label{fig:sankey-experiences}
  \end{subfigure}

  \vspace{.\baselineskip}

  \begin{subfigure}{.5\linewidth}
    \centering
    \includesvg[width=\linewidth]{images/findings/appendix/unfolded-themes-interventions.svg}
    \caption{Human experiences and employed technological interventions.}
    \label{fig:sankey-mechanisms}
  \end{subfigure}

  \vspace{.8\baselineskip}

  \begin{subfigure}{.5\linewidth}
    \centering
    \includesvg[width=\linewidth]{images/findings/appendix/unfolded-interventions-mechanisms.svg}
    \caption{Technological interventions and associated design mechanisms.}
    \label{fig:sankey-interventions}
  \end{subfigure}

  \caption{An unfolded representation of the scoping review findings, showing relationships between human experiences, design mechanisms, and technological interventions.}
  \label{fig:appendix-findings-overview}
\end{widefigure}
